\clearpage\SourceLeaf{111}%
\pdfbookmark[1]{On the Limits of the Geocentric Positions of the Planets}{grenzen1804}
\begin{center}
{\small ON THE LIMITS OF THE GEOCENTRIC POSITIONS OF THE PLANETS.\par}
\medskip\hrule\medskip
{\small Monatliche Correspondenz zur Beförderung der Erd- und Himmels-Kunde,\par
edited by Baron \textsc{von Zach}. August 1804.\par}
\medskip\hrule
\end{center}\medskip
Seen from the Sun, the motion of every planet always appears to follow
one and the same great circle in the sky of the fixed stars, provided
the small perturbations caused by other celestial bodies are disregarded.
It would appear the same from the Earth if the plane of its orbit
coincided with the ecliptic. But if these two planes are inclined to
each other, all the possible geocentric positions of the planet no
longer lie along a line on the celestial sphere, extending in one
dimension as in the former case. Instead they fill an area, a
\textit{zone}, which divides all the rest of the sky, where the planet
can never be seen from the Earth, into two parts, and which may therefore
be called the planet’s \textit{zodiac}. Thus each planet in fact has its
own zodiac, whose boundaries (\textit{limits}) can be determined with
complete precision, provided its orbit and the Earth’s orbit are treated
as conic sections with unchanging elements. The exact determination of
these boundaries is in itself an interesting analytical problem. Its
application, especially to the two new planets, whose zodiacs cover a
considerable area, is also of some practical importance. It is well known
that finding and observing these remarkable celestial bodies requires
very precise and detailed star charts, and that even the best we have
\SourceLeaf{112}%
at present are still far from sufficient for the purpose. Thus, unless
one wishes to draw up detailed charts each year of the regions these
planets pass through, one must consider making a dedicated atlas covering
the whole area in which they can appear. The exact determination of its
boundaries is all the more desirable, since in an undertaking already
of considerable scope one would gladly save any work unnecessary for
the purpose. The news will surely be very welcome to all friends of
astronomy that the skilful Lilienthal astronomer \textsc{Harding}, to
whom we already owe several excellent detailed charts, has begun to
undertake this larger work. It will stand out both for its wealth of
detail and for its exceptionally careful critical scrutiny, based
throughout on firsthand observation.

Through the centre of the Sun we pass three mutually perpendicular
planes, otherwise chosen arbitrarily. We call the perpendicular
distances of the planet under observation from them $x$, $y$, $z$, and
the distances of the Earth $x'$, $y'$, $z'$. We then put
\[
\begin{aligned}
x'-x&=\Delta\cos\alpha\cos\ObservedBeta\\
y'-y&=\Delta\sin\alpha\cos\ObservedBeta\\
z'-z&=\Delta\sin\ObservedBeta
\end{aligned}
\]
so that $\Delta$ is the distance of the planet from the Earth,
$\ObservedBeta$ the inclination of the straight line drawn from the
planet to the Earth to a plane through the Earth parallel to the plane
of the $z$ coordinates, and $\alpha$ the angle between the projection
of that line onto this plane and a plane through the Earth parallel to
the plane of the $y$ coordinates. On the celestial sphere, then, $\alpha$
and $\ObservedBeta$ determine the geocentric position of the planet
relative to the great circles formed by the planes of the $z$, $y$, $x$
coordinates (or, more precisely, by planes drawn through the Earth
parallel to them), just as longitude and latitude determine position
relative to the ecliptic and the equinoctial and solstitial colures.

The nature of the planet’s orbit gives two equations between $x$, $y$,
$z$. We can therefore regard these three variables as functions of
\textit{one} variable, which we denote by $t$ and otherwise leave
unspecified. Likewise let $x'$, $y'$, $z'$ be functions of the variable
\SourceLeaf{113}%
$t'$. Thus $\alpha$ and $\ObservedBeta$ are functions of both $t$, $t'$.
We may specify them by the differential equations
\[
\mathrm{d}\alpha=p\mathrm{d}t+p'\mathrm{d}t',\qquad \mathrm{d}\ObservedBeta=q\mathrm{d}t+q'\mathrm{d}t'
\]

On this assumption, it is clear that if we let $t$, $t'$ vary together
in such a way that $\mathrm{d}t:\mathrm{d}t'=-p':p$, then $\alpha$ stays
unchanged while $\ObservedBeta$ increases or decreases until it has
reached a greatest or least value. This clearly happens when
$pq'-qp'=0$. Now the combination of \textit{all} values of $t$ and $t'$
gives all possible geocentric positions of the planet. Of all the
combinations giving the same $\alpha$, the one for which $\ObservedBeta$
is greatest or least must occur if the geocentric position is to fall
on the limits of the planet’s zodiac. It follows that these limits are
determined by the condition $pq'-qp'=0$.

Differentiating the equations above gives
\[
\begin{aligned}
\mathrm{d}x'-\mathrm{d}x&=\cos\alpha\cos\ObservedBeta\,\mathrm{d}\Delta-\Delta\sin\alpha\cos\ObservedBeta\,\mathrm{d}\alpha-\Delta\cos\alpha\sin\ObservedBeta\,\mathrm{d}\ObservedBeta\\
\mathrm{d}y'-\mathrm{d}y&=\sin\alpha\cos\ObservedBeta\,\mathrm{d}\Delta+\Delta\cos\alpha\cos\ObservedBeta\,\mathrm{d}\alpha-\Delta\sin\alpha\sin\ObservedBeta\,\mathrm{d}\ObservedBeta\\
\mathrm{d}z'-\mathrm{d}z&=\sin\ObservedBeta\,\mathrm{d}\Delta+\Delta\cos\ObservedBeta\,\mathrm{d}\ObservedBeta
\end{aligned}
\]

Combining this with those equations, we readily obtain
\[
\begin{gathered}
-\sin\alpha(\mathrm{d}x'-\mathrm{d}x)+\cos\alpha(\mathrm{d}y'-\mathrm{d}y)=\Delta\cos\ObservedBeta\,\mathrm{d}\alpha\\
\begin{aligned}
-\cos\alpha\sin\ObservedBeta(\mathrm{d}x'-\mathrm{d}x)-\sin\alpha\sin\ObservedBeta(\mathrm{d}y'-\mathrm{d}y)&\\
{}+\cos\ObservedBeta(\mathrm{d}z'-\mathrm{d}z)&=\Delta\mathrm{d}\ObservedBeta
\end{aligned}
\end{gathered}
\]

The partial differentials thus give
\[
\begin{aligned}
\Delta\cos\ObservedBeta\,p\mathrm{d}t&=\sin\alpha\,\mathrm{d}x-\cos\alpha\,\mathrm{d}y\\
\Delta\cos\ObservedBeta\,p'\mathrm{d}t'&=-\sin\alpha\,\mathrm{d}x'+\cos\alpha\,\mathrm{d}y'\\
\Delta q\mathrm{d}t&=\cos\alpha\sin\ObservedBeta\,\mathrm{d}x+\sin\alpha\sin\ObservedBeta\,\mathrm{d}y-\cos\ObservedBeta\,\mathrm{d}z\\
\Delta q'\mathrm{d}t'&=-\cos\alpha\sin\ObservedBeta\,\mathrm{d}x'-\sin\alpha\sin\ObservedBeta\,\mathrm{d}y'+\cos\ObservedBeta\,\mathrm{d}z'
\end{aligned}
\]

Substituting these values of $p$, $p'$, $q$, $q'$ into the condition
$pq'=p'q$ and carrying out the requisite reductions gives
\[
\begin{aligned}
\cos\alpha\cos\ObservedBeta(\mathrm{d}y'\mathrm{d}z-\mathrm{d}y\,\mathrm{d}z')&+\sin\alpha\cos\ObservedBeta(\mathrm{d}z'\mathrm{d}x-\mathrm{d}z\,\mathrm{d}x')\\
&+\sin\ObservedBeta(\mathrm{d}x'\mathrm{d}y-\mathrm{d}x\,\mathrm{d}y')=0
\end{aligned}
\]
\SourceLeaf{114}%
or, on multiplying by $\Delta$,
\[
\begin{aligned}
(x'-x)(\mathrm{d}y'\mathrm{d}z-\mathrm{d}y\,\mathrm{d}z')&+(y'-y)(\mathrm{d}z'\mathrm{d}x-\mathrm{d}z\,\mathrm{d}x')\\
&+(z'-z)(\mathrm{d}x'\mathrm{d}y-\mathrm{d}x\,\mathrm{d}y')=0
\end{aligned}
\]
This equation is better expressed in the following form.
\[
\left.\begin{aligned}
\mathrm{d}x(y'\mathrm{d}z'-z'\mathrm{d}y')&+\mathrm{d}x'(y\,\mathrm{d}z-z\,\mathrm{d}y)\\
{}+\mathrm{d}y(z'\mathrm{d}x'-x'\mathrm{d}z')&+\mathrm{d}y'(z\,\mathrm{d}x-x\,\mathrm{d}z)\\
{}+\mathrm{d}z(x'\mathrm{d}y'-y'\mathrm{d}x')&+\mathrm{d}z'(x\,\mathrm{d}y-y\,\mathrm{d}x)
\end{aligned}\right\}=0
\]

This equation gives in general the relation between the positions of the
Earth and the planet for which the planet’s geocentric position falls on
the limits. One need only substitute for $x$, $y$, $z$ their values in
terms of $t$, and for $x'$, $y'$, $z'$ their values in terms of $t'$, as
determined by the orbits, to obtain a finite equation between $t$ and $t'$.
It seemed worth deriving this equation by a general analysis. It is, however,
easy to show that it is also the condition \textit{that the tangents at the
positions of the Earth and the planet lie in one plane}, and to derive
this very condition from the requirements of our problem. For brevity’s
sake, however, we shall dwell on this no longer here.

For the quantities $t$, $t'$ that we have so far left unspecified, it
is most convenient to take the heliocentric angular distances of the
planet and the Earth in their orbits from the common line of nodes
(specifically, from the ascending node of the planet’s orbit on the
Earth’s orbit). If we now denote the distances of the planet and the
Earth from the Sun by $r$, $r'$, their coordinates can be expressed
as follows:
\[
\begin{aligned}
x&=r\sin a\sin(t+A)\\
y&=r\sin b\sin(t+B)\\
z&=r\sin c\sin(t+C)\\
x'&=r'\sin a'\sin(t'+A')\\
y'&=r'\sin b'\sin(t'+B')\\
z'&=r'\sin c'\sin(t'+C')
\end{aligned}
\]

This, and the meaning of the constants $a$, $A$, and so on, are
readily understood by generalizing the investigation presented in the
May 1804 \textit{Mon. Corr.} [p.\,98 of this volume]. Now let $k$, $k'$
\SourceLeaf{115}%
be the semi-parameters of the conic sections described by the planet
and the Earth, $e$, $e'$ their eccentricities, and $g$, $g'$ the angular
distances of their aphelia from the line of nodes. Then
\[
r=\frac{k}{1-e\cos(t-g)},\qquad r'=\frac{k'}{1-e'\cos(t'-g')}
\]

Carrying out the calculation then gives
\[
\mathrm{d}x=\frac{k\sin a\,\mathrm{d}t}{(1-e\cos(t-g))^2}\times(\cos(t+A)-e\cos(g+A))
\]

The values of $\mathrm{d}y$, $\mathrm{d}z$ have a similar form. To
obtain them one need only replace $a$, $A$ by $b$, $B$ or by $c$, $C$.
The values of $\mathrm{d}x'$, $\mathrm{d}y'$, $\mathrm{d}z'$ are obtained
from those of $\mathrm{d}x$, $\mathrm{d}y$, $\mathrm{d}z$ by replacing the
quantities referring to the planet with their counterparts for the Earth.

It is more convenient to work out $y\,\mathrm{d}z-z\,\mathrm{d}y$
from the values of $y$, $z$ before substituting the value of $r$.
One thus obtains
\[
y\,\mathrm{d}z-z\,\mathrm{d}y=rr\sin b\sin c\sin(B-C)\mathrm{d}t=rr\cos a\,\mathrm{d}t
\]
(see the cited paper [p.\,100 of this volume]), and likewise
\[
\begin{aligned}
z\,\mathrm{d}x-x\,\mathrm{d}z&=rr\cos b\,\mathrm{d}t\\
x\,\mathrm{d}y-y\,\mathrm{d}x&=rr\cos c\,\mathrm{d}t
\end{aligned}
\]

Very similar values are obtained for the three corresponding
expressions for the Earth.

Substituting all these values and carrying out the requisite
reductions gives the following form of the condition above:
\[
\begin{aligned}
&k'\cos a'\sin a\,(\cos(t+A)-e\cos(g+A))\\
&{}+k'\cos b'\sin b\,(\cos(t+B)-e\cos(g+B))\\
&{}+k'\cos c'\sin c\,(\cos(t+C)-e\cos(g+C))\\
&{}+k\cos a\sin a'\,(\cos(t'+A')-e'\cos(g'+A'))\\
&{}+k\cos b\sin b'\,(\cos(t'+B')-e'\cos(g'+B'))\\
&{}+k\cos c\sin c'\,(\cos(t'+C')-e'\cos(g'+C'))=0
\end{aligned}
\]

By suitable reductions, if we denote the inclination of the planet’s
orbit to the Earth’s by $i$, the first three terms of this equation
can be reduced to $k'\sin i\,(\cos t-e\cos g)$ and the last three to
$-k\sin i\,(\cos t'-e'\cos g')$.
\SourceLeaf{116}%
These reductions are in themselves easy, though somewhat lengthy. We
can all the more readily spare ourselves the trouble of developing
them here, since we can reach the same result much more conveniently
by choosing suitably the three fundamental planes to which the
coordinates refer, which we have so far left unspecified. We take
the ecliptic for the plane of the $z$ coordinates and choose the
planes of the $x$, $y$ coordinates so that the pole of the first
lies at the ascending node of the planet’s orbit, and that of the
second lies $90^0$ farther forward in the ecliptic. We take the
coordinates $x$, $y$, $z$ as positive on the side of \textit{these}
poles and on the north side of the ecliptic. It is easy to see that
the following relations hold under these assumptions:
\[
\begin{array}{llll}
a=90^0,&b=90^0+i,&c=i,&C=0\\
a'=90^0,&b'=90^0.&c'=0,&B'=0
\end{array}
\]
The equation above therefore becomes
\[
k'\sin i\,(\cos t-e\cos g)-k\sin i\,(\cos t'-e'\cos g')=0
\]
or
\[
k'(\cos t-e\cos g)=k(\cos t'-e'\cos g')
\]

It is easy to show from the theory of conic sections that
$\frac{k}{\cos t-e\cos g}$ and $\frac{k'}{\cos t'-e'\cos g'}$ express the
distances from the Sun to the points where the tangents at the
positions of the planet and the Earth meet the line of nodes of
the two orbits. The equation just obtained thus says that the two
tangents meet the line of nodes at one and the same point, in
agreement with the condition mentioned above that they lie in one
and the same plane.

There are three cases to distinguish in the position of the planet’s
orbit relative to the Earth’s. Either the planet’s orbit encloses the
Earth’s, or the Earth’s encloses the planet’s, or the two enclose one
another like links of a chain. The first case occurs when the planet
is farther from the Sun than the Earth on both sides along the line
of nodes. The second occurs when the Earth is farther from the Sun
than the planet on both sides. The third occurs when the planet is
farther from the Sun on one side and the Earth on the other. None of
the planets known so far is situated relative to the Earth or to
another planet in the way the third case requires.
\SourceLeaf{117}%
Comets of this kind, however, are plentiful. The analytical condition
for the first case is that $k-k'$ be positive and, irrespective of
sign, greater than $k'e\cos g-ke'\cos g'$. In the second case $k'-k$
must have these properties. In the third case $k-k'$ or $k'-k$ must,
irrespective of sign, be less than $k'e\cos g-ke'\cos g'$.

In the first of these three cases, the equation above gives, for any
value of $t$ from $0^0$ to $360^0$, two values of $t'$. In the
second case each value of $t'$ gives two of $t$. One value of $t'$
(in the first case, or of $t$ in the second) will always lie between
$0$ and $180$, the other between $180$ and $360^0$, or rather within
still narrower limits that are easy to determine more precisely.
In the first case, then, two heliocentric positions of the Earth
correspond to every heliocentric position of the planet. The converse
fails, since the Earth can see the planet at its limits only in two
separate parts of the Earth’s orbit, one below or south of the plane
of the planet’s orbit, the other above or north of it. The planet
appears at its northern limit only when $t'$ takes one of its possible
values between $0$ and $180$, and at its southern limit when it takes
one between $180$ and $360^0$. Likewise, in the second case, two
heliocentric positions of the planet correspond to every heliocentric
position of the Earth. The converse fails here too, since the planet
can appear to the Earth at its limits only in two separate parts of
its orbit, one north of the ecliptic and the other south of it. It
appears at the northern limit for values between $0$ and $180$, and
at the southern limit for values between $180$ and $360^0$.

In the third case, however, both $t$ and $t'$ can take values only
within certain limits. In other words, the Earth and the planet
must each be in a particular part of its orbit if the condition
above is to hold.

It follows that the geocentric positions of the planet arising from
all combinations of heliocentric positions of the planet and the
Earth that satisfy the equation above form, in the first and second
cases, two separate curves on the celestial sphere, each returning
into itself. Between them lies the great circle representing the
plane of the planet’s orbit in the first case, and the ecliptic in the
second case.
\SourceLeaf{118}%
In the third case, however, those geocentric positions form only
\textit{one} curve returning into itself, as a closer look at the
case readily shows.

According to the foregoing investigations, the planet’s zodiac can
have no boundaries other than these curves. It therefore seems natural
to conclude that in the first two cases the planet’s zodiac is the
zone between the two curves, and in the third case one of the two
regions into which the curve divides the entire surface of the sphere.
Yet this conclusion would sometimes fail in the first two cases
and would always fail in the third. Here, as in many other uses of
analysis where a similar precaution is sometimes overlooked, we
must remember that our reasoning rests entirely on the assumption that
the planet’s zodiac really is bounded and that the planet cannot
appear from the Earth at every point in the sky. In the third case,
however, that assumption fails, as can already be shown on grounds
of \textit{geometry of position}. The curve we found can therefore
provide no boundary for the planet’s zodiac in this case, since its
zodiac occupies the \textit{whole} sky. In the first and second cases,
there will indeed always be places in the sky on at least one side
of the zone we found where the planet can never appear, so one of the
curves is certainly a boundary. Yet it can happen that such excluded
places occur on \textit{only} one side. The other curve then gives
no boundary. The region of the sky it cuts off from the zone belongs
to the zodiac in its entirety, just as the zone itself does. This
is hardly the place to pursue the investigation in full and to
explain what those curves actually mean in cases where they give
no boundaries. We can all the more readily content ourselves here
with drawing the attention of those interested in analysis to these
seemingly paradoxical phenomena, since it is easy to show that none
of the planets known so far, which are our immediate concern here,
ever appears from the Earth at either pole of the ecliptic. The
exceptions mentioned therefore cannot occur for these planets.
Their zodiacs must therefore be genuine \textit{zones}, and the two
curves found must be their boundaries.

We shall now apply our results to \textit{Pallas} and \textit{Ceres},
partly for further illustration and partly for practical use.
\SourceLeaf{119}%
We shall mark out the limits of their zodiacs so that they can be
entered on star charts. The perturbations will indeed require some
extension of these limits, and the changes the elements will undergo
in the future will call for some alterations. In practice, however,
these changes will be insignificant. If after some years observations
make them perceptible, they will be particularly interesting for that
very reason, since they will then show conspicuously how the elliptic
elements used here gradually cease to suffice.

For \textit{Pallas} we take the following values from the latest
elements for the beginning of 1803 (\textit{Mon. Corr.}, March 1804).
\[
\begin{array}{rl@{\qquad}rl}
e&=0.2457396&\varphi&=14^0\;13'\;31''966\\
k&=2.602122&&\\
g&=128^0\;49'\;20''7&&\\
\AscendingNode&=172\;28\;13.7&i&=34\;38\;1.1
\end{array}
\]

For the Earth, on the other hand,
\[
\begin{array}{rl@{\qquad}rl}
e'&=0.016792&\varphi&=0^0\;57'\;43''802\\
k'&=0.999718&&\\
g'&=107^0\;5'\;39''&&\\
\mathrm{Aph.}&=279\;33\;52.7&&
\end{array}
\]

Substituting these values gives the following form of our equation above.
\[
\cos t'=0.384193\cos t+0.0542514
\]

It follows that the two extreme values of $\cos t'$ are $+0.438444$
and $-0.329942$. Thus all possible values of $t'$ lie either between
$63^0\;59'\;43''$ and $109^0\;15'\;55''$ or between $250^0\;44'\;5''$
and $296^0\;0'\;17''$. Hence \textit{Pallas} can appear to the Earth
at its limits only when the Earth’s heliocentric longitude is
between $236^0\;28'$ and $281^0\;44'$ or between $63^0\;12'$ and
$180^0\;29'$, that is, roughly from 18\,May to 4\,July and from
26\,November to 9\,January. In the first part of its orbit the Earth
is south of the plane of Pallas’s orbit, and in the other it is
north of it. Thus \textit{Pallas} will appear to the Earth at the
northern limit of its zodiac in the first part and at the southern
limit in the other. It is also easy to show that \textit{Pallas}
\SourceLeaf{120}%
must touch the northern limit once and the southern limit once each
year at the specified times. --- To obtain enough points from both
limits, we shall take for $t$ in succession all values from $0$ to
$360^0$, from one multiple of $10$ degrees to the next in steps of
$10$ degrees. By combining each with the two corresponding values of $t'$ determined from the formula above, we shall obtain
the corresponding geocentric positions directly in right ascension
and declination. For this purpose we can use the method explained
in the paper mentioned above and the constants already calculated
there. The following table gives the results of these calculations.

\begin{center}
\begingroup\small\setlength{\tabcolsep}{5pt}
\begin{tabular}{r||rr|rrl||rr|rrl}
&\multicolumn{5}{c||}{Northern limit}&\multicolumn{5}{c}{Southern limit}\\
$t$&\multicolumn{2}{c|}{Right asc.}&\multicolumn{3}{c||}{Declination}&\multicolumn{2}{c|}{Right asc.}&\multicolumn{3}{c}{Declination.}\\\hline
$0^0$&$148^0$&$4'$&$12^0$&$56'$&N.&$196^0$&$54'$&$7^0$&$12'$&S.\\
$10$&$159$&$30$&$15$&$56$&&$205$&$54$&$4$&$26$&\\
$20$&$171$&$54$&$18$&$55$&&$214$&$15$&$2$&$11$&\\
$30$&$185$&$6$&$21$&$38$&&$222$&$3$&$0$&$23$&\\
$40$&$198$&$48$&$23$&$52$&&$229$&$27$&$1$&$1$&N.\\
$50$&$212$&$41$&$25$&$26$&&$236$&$32$&$2$&$4$&\\
$60$&$226$&$28$&$26$&$16$&&$243$&$26$&$2$&$48$&\\
$70$&$239$&$59$&$26$&$23$&&$250$&$16$&$3$&$16$&\\
$80$&$253$&$8$&$25$&$53$&&$257$&$6$&$3$&$28$&\\
$90$&$265$&$54$&$24$&$49$&&$264$&$1$&$3$&$25$&\\
$100$&$278$&$18$&$23$&$19$&&$271$&$4$&$3$&$5$&\\
$110$&$290$&$25$&$21$&$27$&&$278$&$18$&$2$&$30$&\\
$120$&$302$&$19$&$19$&$17$&&$285$&$44$&$1$&$37$&\\
$130$&$314$&$4$&$16$&$53$&&$293$&$23$&$0$&$27$&\\
$140$&$325$&$43$&$14$&$22$&&$301$&$16$&$1$&$2$&S.\\
$150$&$337$&$16$&$11$&$47$&&$309$&$24$&$2$&$50$&\\
$160$&$348$&$42$&$9$&$14$&&$317$&$49$&$4$&$57$&\\
$170$&$359$&$55$&$6$&$49$&&$326$&$34$&$7$&$23$&\\
$180$&$10$&$49$&$4$&$40$&&$335$&$41$&$10$&$8$&\\
$190$&$21$&$15$&$2$&$49$&&$345$&$12$&$13$&$10$&\\
$200$&$31$&$5$&$1$&$21$&&$355$&$12$&$16$&$24$&\\
$210$&$40$&$14$&$0$&$14$&&$5$&$44$&$19$&$53$&\\
$220$&$48$&$39$&$0$&$31$&S.&$16$&$53$&$23$&$6$&\\
$230$&$56$&$23$&$0$&$59$&&$28$&$41$&$26$&$16$&\\
$240$&$63$&$28$&$1$&$10$&&$41$&$12$&$29$&$5$&\\
$250$&$70$&$2$&$1$&$9$&&$54$&$25$&$31$&$21$&\\
$260$&$76$&$12$&$0$&$57$&&$68$&$16$&$32$&$52$&\\
$270$&$82$&$5$&$0$&$35$&&$82$&$36$&$33$&$29$&\\
$280$&$87$&$51$&$0$&$4$&&$97$&$12$&$33$&$5$&\\
$290$&$93$&$38$&$0$&$37$&N.&$111$&$47$&$31$&$39$&\\
$300$&$99$&$35$&$1$&$30$&&$126$&$6$&$29$&$14$&\\
$310$&$105$&$53$&$2$&$36$&&$139$&$52$&$26$&$0$&\\
$320$&$112$&$40$&$3$&$59$&&$152$&$57$&$22$&$12$&\\
$330$&$120$&$8$&$5$&$40$&&$165$&$12$&$18$&$9$&\\
$340$&$128$&$26$&$7$&$44$&&$176$&$37$&$14$&$10$&\\
$350$&$137$&$43$&$10$&$10$&&$187$&$10$&$10$&$28$&\\
$360$&$148$&$4$&$12$&$56$&&$196$&$54$&$7$&$12$&
\end{tabular}
\endgroup
\end{center}
\SourceLeaf{121}%
For greater convenience, suitable interpolation methods were used
to insert between these $72$ points the following $144$ points. The
right ascensions increase in steps of $5$ degrees, from one multiple
of $5$ degrees to the next.

\begin{center}
\begingroup\small\setlength{\tabcolsep}{3pt}
\begin{tabular}{r|rrl|rrl||r|rrl|rrl}
\multicolumn{14}{c}{\textit{Zodiac of Pallas}}\\[4pt]
&\multicolumn{6}{c||}{Declination}&&\multicolumn{6}{c}{Declination}\\
\shortstack{Right\\asc.}&\multicolumn{3}{c|}{\shortstack{at northern\\limit}}&\multicolumn{3}{c||}{\shortstack{at southern\\limit}}&\shortstack{Right\\asc.}&\multicolumn{3}{c|}{\shortstack{at northern\\limit}}&\multicolumn{3}{c}{\shortstack{at southern\\limit}}\\\hline
$0^0$&$6^0$&$48'$&N.&$17^0$&$56'$&S.&$180^0$&$20^0$&$38'$&N.&$12^0$&$58'$&S.\\
$5$&$5$&$47$&&$19$&$32$&&$185$&$21$&$37$&&$11$&$12$&\\
$10$&$4$&$49$&&$21$&$4$&&$190$&$22$&$31$&&$9$&$29$&\\
$15$&$3$&$54$&&$22$&$33$&&$195$&$23$&$19$&&$7$&$49$&\\
$20$&$3$&$2$&&$23$&$59$&&$200$&$24$&$2$&&$6$&$13$&\\
$25$&$2$&$13$&&$25$&$20$&&$205$&$24$&$39$&&$4$&$42$&\\
$30$&$1$&$29$&&$26$&$36$&&$210$&$25$&$11$&&$3$&$17$&\\
$35$&$0$&$50$&&$27$&$46$&&$215$&$25$&$37$&&$2$&$0$&\\
$40$&$0$&$16$&&$28$&$51$&&$220$&$25$&$58$&&$0$&$49$&\\
$45$&$0$&$14$&S.&$29$&$49$&&$225$&$26$&$12$&&$0$&$13$&N.\\
$50$&$0$&$37$&&$30$&$41$&&$230$&$26$&$22$&&$1$&$7$&\\
$55$&$0$&$55$&&$31$&$26$&&$235$&$26$&$25$&&$1$&$52$&\\
$60$&$1$&$6$&&$32$&$4$&&$240$&$26$&$23$&&$2$&$28$&\\
$65$&$1$&$11$&&$32$&$36$&&$245$&$26$&$16$&&$2$&$56$&\\
$70$&$1$&$9$&&$33$&$0$&&$250$&$26$&$3$&&$3$&$15$&\\
$75$&$1$&$0$&&$33$&$17$&&$255$&$25$&$45$&&$3$&$26$&\\
$80$&$0$&$44$&&$33$&$27$&&$260$&$25$&$23$&&$3$&$28$&\\
$85$&$0$&$21$&&$33$&$30$&&$265$&$24$&$55$&&$3$&$23$&\\
$90$&$0$&$10$&N.&$33$&$25$&&$270$&$24$&$22$&&$3$&$9$&\\
$95$&$0$&$49$&&$33$&$13$&&$275$&$23$&$46$&&$2$&$48$&\\
$100$&$1$&$34$&&$32$&$54$&&$280$&$23$&$5$&&$2$&$19$&\\
$105$&$2$&$27$&&$32$&$27$&&$285$&$22$&$20$&&$1$&$43$&\\
$110$&$3$&$25$&&$31$&$53$&&$290$&$21$&$31$&&$1$&$0$&\\
$115$&$4$&$29$&&$31$&$11$&&$295$&$20$&$39$&&$0$&$10$&\\
$120$&$5$&$38$&&$30$&$23$&&$300$&$19$&$43$&&$0$&$46$&S.\\
$125$&$6$&$52$&&$29$&$27$&&$305$&$18$&$45$&&$1$&$49$&\\
$130$&$8$&$9$&&$28$&$24$&&$310$&$17$&$44$&&$2$&$58$&\\
$135$&$9$&$27$&&$27$&$14$&&$315$&$16$&$42$&&$4$&$12$&\\
$140$&$10$&$47$&&$25$&$57$&&$320$&$15$&$37$&&$5$&$32$&\\
$145$&$12$&$7$&&$24$&$35$&&$325$&$14$&$31$&&$6$&$56$&\\
$150$&$13$&$27$&&$23$&$7$&&$330$&$13$&$24$&&$8$&$24$&\\
$155$&$14$&$46$&&$21$&$33$&&$335$&$12$&$17$&&$9$&$56$&\\
$160$&$16$&$4$&&$19$&$55$&&$340$&$11$&$10$&&$11$&$30$&\\
$165$&$17$&$19$&&$18$&$14$&&$345$&$10$&$3$&&$13$&$6$&\\
$170$&$18$&$29$&&$16$&$30$&&$350$&$8$&$57$&&$14$&$43$&\\
$175$&$19$&$36$&&$14$&$44$&&$355$&$7$&$52$&&$16$&$20$&\\
$180$&$20$&$38$&&$12$&$58$&&$360$&$6$&$48$&&$17$&$56$&
\end{tabular}
\endgroup
\end{center}

For \textit{Ceres}, according to the latest elements for 1803, we have
\[
\begin{aligned}
e&=0.0788941.\\
k&=2.750681\\
g&=245^0\;34'\;56''
\end{aligned}
\]
\SourceLeaf{122}
For the Earth, $e'$ and $k'$ are as above, and $g'=198^0\;35'\;31''$.
This gives the following condition.
\[
\cos t'=.\,0.363444\cos t-0.004063
\]

The values of $t'$ therefore range from $68^0\;56'\;16''$ to
$111^0\;33'\;43''$ and from $248^0\;26'\;17''$ to $291^0\;3'\;44''$.
The Earth’s heliocentric longitude thus ranges from $149^0\;55'$
to $192^0\;32'$ and from $329^0\;25'$ to $12^0\;2'$. The Earth can
therefore see \textit{Ceres} at its northern limits only roughly
from 19\,February to 3\,April, and at its southern limits from
23\,August to 6\,October. This formula was used, just as before
for \textit{Pallas}, to calculate 36 points on each limit of the
zodiac of \textit{Ceres} in right ascension and declination. In
connection with the paper already mentioned several times, the
following values were used for the constants $a$, $A$, and so on.
\[
\begin{array}{rcrrr}
a&=&79^0&30'&5''\\
b&=&114&42&11\\
c&=&27&7&24\\
A&=&170&49&4\\
B&=&85&42&29\\
C&=&59&36&33
\end{array}
\]
Thus, if $v$ denotes the true anomaly of \textit{Ceres}, the three
coordinates are given by the following formulas.
\[
\begin{aligned}
x&=\frac{\alpha\sin(v+56^0\;24'\;0'')}{1-e\cos v}\\[8pt]
y&=\frac{\ObservedBeta\sin(v+331^0\;17'\;25'')}{1-e\cos v}\\[8pt]
z&=\frac{\gamma\sin(v+305^0\;11'\;29'')}{1-e\cos v}
\end{aligned}
\]
where
\[
\begin{aligned}
\log\alpha&=\log k\sin a=0.432108\\
\log\ObservedBeta&=\log k\sin b=0.397758\\
\log\gamma&=\log k\sin c=0.098319
\end{aligned}
\]
\SourceLeaf{123}

In place of these $72$ points, we content ourselves here with giving
in the following table only the $144$ points inserted between them
in much the same way as for \textit{Pallas}.

\begin{center}
\begingroup\small\setlength{\tabcolsep}{3pt}
\begin{tabular}{r|rrl|rrl||r|rrl|rrl}
\multicolumn{14}{c}{\textit{Zodiac of Ceres}}\\[4pt]
&\multicolumn{6}{c||}{Declination}&&\multicolumn{6}{c}{Declination}\\
\shortstack{Right\\asc.}&\multicolumn{3}{c|}{\shortstack{at northern\\limit}}&\multicolumn{3}{c||}{\shortstack{at southern\\limit}}&\shortstack{Right\\asc.}&\multicolumn{3}{c|}{\shortstack{at northern\\limit}}&\multicolumn{3}{c}{\shortstack{at southern\\limit}}\\\hline
$0^0$&$8^0$&$37'$&S.&$17^0$&$29'$&S.&$180^0$&$18^0$&$49'$&N.&$8^0$&$17'$&N.\\
$5$&$6$&$15$&&$15$&$13$&&$185$&$16$&$31$&&$5$&$56$&\\
$10$&$3$&$45$&&$12$&$51$&&$190$&$14$&$5$&&$3$&$28$&\\
$15$&$1$&$11$&&$10$&$26$&&$195$&$11$&$33$&&$0$&$55$&\\
$20$&$1$&$27$&N.&$7$&$56$&&$200$&$8$&$57$&&$1$&$41$&S.\\
$25$&$4$&$7$&&$5$&$18$&&$205$&$6$&$16$&&$4$&$19$&\\
$30$&$6$&$48$&&$2$&$38$&&$210$&$3$&$34$&&$6$&$56$&\\
$35$&$9$&$26$&&$0$&$3$&&$215$&$0$&$52$&&$9$&$32$&\\
$40$&$12$&$1$&&$2$&$30$&N.&$220$&$1$&$47$&S.&$12$&$3$&\\
$45$&$14$&$31$&&$4$&$58$&&$225$&$4$&$21$&&$14$&$29$&\\
$50$&$16$&$54$&&$7$&$20$&&$230$&$6$&$50$&&$16$&$48$&\\
$55$&$19$&$9$&&$9$&$35$&&$235$&$9$&$12$&&$18$&$59$&\\
$60$&$21$&$16$&&$11$&$42$&&$240$&$11$&$25$&&$21$&$1$&\\[3pt]
$65$&$23$&$12$&&$13$&$39$&&$245$&$13$&$28$&&$22$&$53$&\\
$70$&$24$&$58$&&$15$&$27$&&$250$&$15$&$21$&&$24$&$35$&\\
$75$&$26$&$34$&&$17$&$4$&&$255$&$17$&$3$&&$26$&$6$&\\
$80$&$27$&$59$&&$18$&$30$&&$260$&$18$&$34$&&$27$&$26$&\\
$85$&$29$&$13$&&$19$&$45$&&$265$&$19$&$53$&&$28$&$36$&\\
$90$&$30$&$16$&&$20$&$49$&&$270$&$21$&$1$&&$29$&$35$&\\[3pt]
$95$&$31$&$8$&&$21$&$42$&&$275$&$21$&$56$&&$30$&$24$&\\
$100$&$31$&$50$&&$22$&$24$&&$280$&$22$&$40$&&$31$&$1$&\\
$105$&$32$&$21$&&$22$&$55$&&$285$&$23$&$13$&&$31$&$29$&\\
$110$&$32$&$42$&&$23$&$14$&&$290$&$23$&$34$&&$31$&$46$&\\
$115$&$32$&$51$&&$23$&$22$&&$295$&$23$&$43$&&$31$&$52$&\\
$120$&$32$&$50$&&$23$&$19$&&$300$&$23$&$41$&&$31$&$48$&\\[3pt]
$125$&$32$&$40$&&$23$&$5$&&$305$&$23$&$28$&&$31$&$34$&\\
$130$&$32$&$18$&&$22$&$40$&&$310$&$23$&$4$&&$31$&$9$&\\
$135$&$31$&$46$&&$22$&$3$&&$315$&$22$&$27$&&$30$&$34$&\\
$140$&$31$&$3$&&$21$&$16$&&$320$&$21$&$40$&&$29$&$49$&\\
$145$&$30$&$9$&&$20$&$17$&&$325$&$20$&$40$&&$28$&$53$&\\
$150$&$29$&$4$&&$19$&$6$&&$330$&$19$&$30$&&$27$&$46$&\\[3pt]
$155$&$27$&$49$&&$17$&$45$&&$335$&$18$&$8$&&$26$&$29$&\\
$160$&$26$&$22$&&$16$&$12$&&$340$&$16$&$35$&&$25$&$1$&\\
$165$&$24$&$44$&&$14$&$28$&&$345$&$14$&$51$&&$23$&$23$&\\
$170$&$22$&$56$&&$12$&$34$&&$350$&$12$&$55$&&$21$&$34$&\\
$175$&$20$&$57$&&$10$&$30$&&$355$&$10$&$51$&&$19$&$36$&\\
$180$&$18$&$49$&&$8$&$17$&&$360$&$8$&$37$&&$17$&$29$&\\
\end{tabular}
\endgroup
\end{center}
